\documentclass{article}

% ------------------------------------------------------------
% First publication (survey). Separate from main.tex (the audit, "Bөlinbeitin").
% arxiv.sty is in this directory (kourgeorge/arxiv-style).
% Build:  pdflatex survey.tex && pdflatex survey.tex
% (bibliography is inline via thebibliography, so no bibtex pass needed for the draft)
% ------------------------------------------------------------
\usepackage{arxiv}

\usepackage[utf8]{inputenc}
\usepackage[T1]{fontenc}
\usepackage[english]{babel}
% NOTE: the typeset content is English-only. If Kazakh examples in Cyrillic are added later,
% restore  \usepackage[T2A,T1]{fontenc}  and  \usepackage[russian,english]{babel}
% (requires the texlive-lang-cyrillic package locally; arXiv's TeX Live already ships it).
\usepackage{hyperref}
\usepackage{url}
\usepackage{booktabs}
\usepackage{amsmath}
\usepackage{microtype}
\usepackage{xcolor}

\newcommand{\todo}[1]{\textcolor{red}{[\textbf{TODO}: #1]}}

\title{
  The State of Kazakh NLP: \\
  A Data-Driven Survey from Global Breakthroughs \\
  to an Open Research Agenda
}

\author{
  Altair Zhambyl \\
  Independent Researcher \\
  Almaty, Kazakhstan \\
  \texttt{hey@evgenyastapov.com}
}

% ------------------------------------------------------------
\begin{document}
\maketitle

\begin{abstract}
Kazakh, an agglutinative Turkic language spoken by roughly 13 million people, has
moved in a few years from near-absence in the NLP literature to a small but rapidly
growing research field. Yet existing surveys treat Kazakh only as one of several Central
Asian Turkic languages, leaving newcomers without a field-level map of what has been done
in Kazakh specifically and where the open problems lie. We present the first survey
dedicated to Kazakh, and the first built by a \emph{reproducible}, citation-grounded
pipeline rather than hand-curation: we automatically assemble a corpus of \textbf{222}
Kazakh-focused papers
(2013--2026) from the arXiv and Semantic Scholar APIs and analyse the resulting
citation graph. Three findings organise the field. First, Kazakh NLP exhibits a
\textbf{4--7 year adoption lag} behind global architectural breakthroughs (the
Transformer appears in 2017 but the first large Kazakh LLMs only in 2024--2025).
Second, the field rests on a handful of \emph{dataset hubs}---the Kazakh Speech Corpus
is cited 14 times within the corpus---and on four global pillars (the Transformer, BPE,
BERT, and XLM-R). Third, the field is both young and lopsided: \textbf{47\%} of all
papers (105 of 222) appeared in 2024--2026, while across all years effort concentrates in
speech (110 papers) and machine translation (105), leaving tokenization (21) and
morphology (48) a thin, under-charted frontier. We close with an actionable
research agenda of cited pain points and unverified hypotheses aimed at first-time
contributors. All code, the corpus, and an interactive atlas are released to support
reproduction and extension.
\end{abstract}

\section{Introduction}
Large language models have reorganised natural language processing around a small set
of architectural and training breakthroughs, but the benefits of those breakthroughs
reach the world's languages unevenly. Kazakh---an agglutinative Turkic language with
rich suffixal morphology and vowel harmony---is a representative case of a
\emph{moderately resourced} language: too central to a nation-state to be ignored,
too far from the high-resource core to be well served by multilingual models trained
predominantly on English.

In the span of roughly five years, Kazakh NLP has grown from scattered speech-recognition
experiments into a field with dedicated corpora, named-entity and question-answering
datasets, sentiment resources, and---most recently---several from-scratch and adapted
large language models. This growth has been fast and largely uncoordinated, and no
survey has yet stepped back to ask: \emph{what, exactly, has been built, in what order,
on whose shoulders, and what remains open?}

This paper answers that question with three contributions:
\begin{enumerate}
  \item \textbf{A reproducible, data-driven survey method.} Rather than hand-pick papers,
  we assemble the corpus automatically from public APIs and release the scraping and
  analysis code, so the survey can be re-run and extended as the field grows
  (Section~\ref{sec:method}).
  \item \textbf{A quantitative genealogy of the field.} Using the citation graph, we order
  the work as a citation-derived sequence of breakthroughs, surface the dataset and
  architectural hubs the field depends on, and measure the adoption lag behind global research
  (Sections~\ref{sec:foundations}--\ref{sec:models}).
  \item \textbf{An open research agenda.} We compile cited pain points and unverified
  hypotheses scoped to be tractable for first-time contributors, turning a passive survey
  into a roadmap (Section~\ref{sec:agenda}).
\end{enumerate}

To our knowledge this is the first survey dedicated specifically to Kazakh, and the first
built by a reproducible, citation-grounded pipeline. The closest prior work surveys Kazakh
as one of four Central Asian Turkic languages~\cite{loreslm} or reviews Turkic NLP at
large~\cite{turkicreview}; others treat Kazakh as one of many languages in a broader
cross-lingual study---e.g.\ the 70-language morphological-alignment study of Arnett et
al.~\cite{arnett70}. None provides the field-level, single-language, citation-grounded
genealogy we present.

\section{Methodology}
\label{sec:method}
\paragraph{Corpus construction.}
We query the arXiv API and the Semantic Scholar Graph API (both accessed in June 2026) for
Kazakh-related NLP/ML work. A paper is retained only if (i) its title or abstract mentions
\emph{Kazakh} or \emph{Qazaq}, and (ii) it is relevant to NLP/ML (a keyword filter over
tasks such as language modelling, tokenization, morphology, translation, and speech).
Records are de-duplicated across sources by arXiv id, DOI, or normalised title. The
resulting corpus contains \textbf{222} papers spanning 2013--2026; because of the API
coverage limits noted below, all counts in this survey are best read as lower bounds on the
true literature rather than exact censuses.

\paragraph{Citation graph.}
For each resolvable paper we retrieve its reference list via the Semantic Scholar
\texttt{batch} endpoint and keep edges whose target is also in the corpus
(\emph{intra-corpus} edges) or is one of a small set of global breakthrough papers.
This yields \textbf{312} edges (182 intra-corpus and 130 to global hubs); 205 of 222
papers were resolved and 104 form the connected core. We use in-degree (citations
received within the corpus) as a measure of local influence and a topological view of
the graph to order the field chronologically and logically
(Section~\ref{sec:genealogy}).

\paragraph{Task tagging.}
Each paper is assigned one or more of ten task categories (tokenization, morphology,
language modelling, machine translation, speech, NER/IE, dataset/corpus, evaluation,
classification, embeddings) by a keyword heuristic over title and abstract.

\paragraph{Reproducibility and scope.}
The scrapers, the corpus, the citation graph, and an interactive atlas are released as a
public artifact.\footnote{Code, corpus, citation graph, and interactive atlas:
\url{https://github.com/altairzhambyl/kazakh-nlp-survey}.}
We note two limitations of the method up front (expanded in Section~\ref{sec:limits}):
the Semantic Scholar API was intermittently rate-limited, so coverage of non-arXiv
venues in some categories is incomplete; and category tags are heuristic rather than
manually adjudicated.

\section{Global Foundations}
\label{sec:foundations}
Kazakh NLP is built almost entirely on architectures developed elsewhere. The citation
graph makes the dependence explicit: the most-cited external works, ranked by how many
Kazakh papers reference them, are the Transformer~\cite{transformer} (23 citing papers),
byte-pair encoding for NMT~\cite{bpe} (17), BERT~\cite{bert} (16), XLM-R~\cite{xlmr} (15),
attention-based translation~\cite{bahdanau} (12), and GPT-3~\cite{gpt3} (10). Tokenization
research enters the picture through Sennrich et al.'s subword units~\cite{bpe} and, later,
fertility-based analyses of multilingual tokenizers~\cite{rust}. These four
``pillars''---a contextual architecture (Transformer/BERT), a multilingual encoder
(XLM-R), a subword algorithm (BPE), and a scaling recipe (GPT-3)---account for the bulk of
external influence on the field.

\section{A Genealogy of Kazakh NLP}
\label{sec:genealogy}
Ordering the citation graph topologically yields a reproducible sequence of breakthroughs
that we group into five eras. This order is one of several valid topological sortings; we
break ties by publication year, so the sequence below should be read as a structural
narrative rather than a unique ranking.

\paragraph{Era I --- Speech first (2015--2019).}
Kazakh NLP grew out of speech and translation rather than text understanding. The earliest
hubs are bilingual Kazakh--Russian ASR systems and English--Kazakh NMT efforts that build
directly on attention and BPE~\cite{bahdanau,bpe,transformer}.

\paragraph{Era II --- The foundational corpus (2020).}
The single most influential resource is the crowdsourced Kazakh Speech Corpus
(KSC)~\cite{ksc}, cited 14 times within our corpus---more than any other Kazakh work. It
converts speech recognition from a bespoke effort into a shared benchmark and seeds most
of what follows.

\paragraph{Era III --- The dataset-building years (2021--2022).}
A wave of open resources builds on KSC: KazNERD for named-entity
recognition~\cite{kaznerd} (11 citations), the KazakhTTS~\cite{kazakhtts} and
KazakhTTS2~\cite{kazakhtts2} text-to-speech corpora (7 and 8), and the industrial-scale
KSC2~\cite{ksc2} (7). The field is, structurally, \emph{dataset-driven}: progress is
gated by corpora, and a handful of them carry most downstream work.

\paragraph{Era IV --- LLMs and evaluation (2024).}
Attention turns to large models and to measuring them. ``Do LLMs Speak
Kazakh?''~\cite{dollms} provides the first systematic evaluation of seven models (6
citations); KazQAD~\cite{kazqad}
and KazSAnDRA~\cite{kazsandra} add open-domain QA and sentiment, and KazMMLU~\cite{kazmmlu}
and TUMLU~\cite{tumlu} bring knowledge benchmarks. These works depend jointly on the earlier
datasets and on global encoders~\cite{bert,xlmr}.

\paragraph{Era V --- From-scratch and adapted LLMs (2025--2026).}
The frontier is Kazakh-specific language models: Sherkala~\cite{sherkala}, the from-scratch
SozKZ family~\cite{sozkz}, the tokenizer-free KazByte adapter~\cite{kazbyte}, and the
state-backed Til-Core.\footnote{Til-Core-0.5B model card:
\url{https://huggingface.co/TilQazyna/Til-Core-0.5B}. At the time of writing it has no
accompanying paper or published benchmark results.} The adoption lag is now quantifiable: the Transformer (2017) is
cited throughout, but the first large Kazakh LLMs appear only in 2024--2025---a 4--7 year
delay between a global breakthrough and its Kazakh realisation.

\section{The Landscape by Task}
\label{sec:landscape}
Table~\ref{tab:categories} shows the distribution of the corpus across task categories.
Because tags are multi-label, the figures are tag counts rather than a partition of the
222 papers. The shape is informative: Kazakh NLP is, historically, a
\emph{speech-and-translation} field. Dataset/corpus papers (146) dominate because resource
construction is the binding
constraint; speech (110) and machine translation (105) follow. Crucially, the two
categories most relevant to the language's defining property---its agglutinative
morphology---are the thinnest: morphology (48) and tokenization (21). This is the central
empirical observation of the survey: the field has invested heavily where data could be
collected, and least where the language is hardest. The corpus is also sharply concentrated
in time: 105 of the 222 papers (47\%) carry a 2024--2026 date, against 18 in 2020 and 3 at
the earliest point in 2013, confirming that Kazakh NLP is overwhelmingly a development of
the last three years.

\begin{table}[t]
  \centering
  \begin{tabular}{lr}
    \toprule
    \textbf{Task category} & \textbf{Papers tagged} \\
    \midrule
    Dataset / corpus      & 146 \\
    Speech (ASR / TTS)    & 110 \\
    Machine translation   & 105 \\
    Evaluation / benchmark & 76 \\
    NER / information extraction & 55 \\
    Language modelling / LLM & 50 \\
    Morphology / segmentation & 48 \\
    Classification / sentiment & 26 \\
    Embeddings            & 22 \\
    \textbf{Tokenization} & \textbf{21} \\
    \bottomrule
  \end{tabular}
  \caption{Corpus by task category. Tags are multi-label, so a paper may be counted under
  several categories: the column reports the number of papers carrying each tag and sums to
  659 tag occurrences across the 222 papers, not to 222. Tokenization and
  morphology---most relevant to Kazakh's agglutinative structure---are the least studied.}
  \label{tab:categories}
\end{table}

\section{Models and Tokenizers}
\label{sec:models}
Table~\ref{tab:models} compares the flagship Kazakh-facing models. Two patterns stand out.
First, the field is converging on \emph{Kazakh-specific tokenization}: vocabularies grow
from 32k (KazRoBERTa) to 256k (Til-Core), and approaches diverge from extended BPE
(Sherkala) to from-scratch BPE (SozKZ) to abandoning the tokenizer entirely
(KazByte~\cite{kazbyte}). Second, claims are running ahead of evidence: the state-backed
Til-Core advertises a morpheme-constrained tokenizer (``support for Kazakh morphology'')
but has, at the time of writing, \emph{no published benchmark results}. Whether
morphology-aware tokenization actually helps Kazakh remains an open question---and a
contested one, since prior work on agglutinative languages finds morphological
segmentation does not reliably beat BPE downstream~\cite{saleva}.

\begin{table}[t]
  \centering
  \small
  \begin{tabular}{lllll}
    \toprule
    \textbf{Model} & \textbf{Year} & \textbf{Vocab} & \textbf{Morphology claim} & \textbf{Benchmarks} \\
    \midrule
    Til-Core-0.5B   & 2025 & 256k (morphBPE) & yes (segmenter unreleased) & \textbf{none} \\
    Sherkala-Chat-8B & 2025 & 159k & no (fertility-driven) & yes \\
    SozKZ (50M--600M) & 2026 & 50k (from scratch) & no & partial \\
    KazByte          & 2026 & byte-level (none) & n/a & yes \\
    KazLLM (8B/70B)  & 2024 & 128k (Llama-3.1) & no & yes \\
    KazRoBERTa       & 2022 & 32k & no & partial \\
    \bottomrule
  \end{tabular}
  \caption{Flagship Kazakh-facing models. The ``morphology claim'' column reports authors'
  assertions, not verified results. Entries without an accompanying paper (Til-Core, KazLLM,
  KazRoBERTa) are characterised from their public model cards rather than peer-reviewed work.}
  \label{tab:models}
\end{table}

\section{An Open Research Agenda}
\label{sec:agenda}
We distil the survey into concrete entry points for new contributors, separating
\emph{pain points} (capabilities or resources the field lacks) from \emph{unverified
hypotheses} (claims that are asserted or assumed but not rigorously tested).

\subsection{Pain points}
Representative, tractable gaps include: the high \emph{tokenization tax} on Kazakh, which
motivates KazByte~\cite{kazbyte} and SozKZ~\cite{sozkz} but is not measured downstream;
the absence of a standardised \emph{morpheme-boundary benchmark}; weak support for
Kazakh--Russian \emph{code-switching}, despite its prevalence in real usage; a
\emph{fragmented evaluation landscape} (KazMMLU, TUMLU, KazQAD, Qorgau use different
protocols~\cite{kazmmlu,tumlu,kazqad,qorgau}); near-absent ASR for spontaneous and
children's speech; and the failure of multimodal models on Kazakh Arabic/Latin-script OCR.
Further tractable gaps include sentiment analysis beyond product reviews, intrinsic
embedding benchmarks for Kazakh, and ASR punctuation restoration. The released agenda lists
each with its motivating citation.

\subsection{Unverified hypotheses}
The field carries several load-bearing assumptions that a student could test:
(1) morphology-aware segmentation improves downstream Kazakh tasks over BPE---\emph{contested}
by~\cite{saleva}; (2) byte-level tokenization beats BPE for agglutinative
Kazakh---\emph{unverified}, with KazByte's authors noting validation is ongoing~\cite{kazbyte};
(3) cross-lingual transfer from Turkish beats transfer from Russian---\emph{assumed} but
never ablated; (4) larger tokenizer vocabularies meaningfully improve downstream
quality---\emph{unverified}, with no controlled study; (5) a small from-scratch Kazakh model
beats a large multilingual one at equal inference budget---\emph{unverified}. Each is stated
with a concrete experiment in the released agenda.

\section{Limitations}
\label{sec:limits}
This survey is bounded by its method. Semantic Scholar rate-limiting left non-arXiv
coverage incomplete in several categories, so counts are lower bounds rather than exact
censuses. Task tags are heuristic and were not manually adjudicated. The citation graph
captures only references resolvable through Semantic Scholar, so the connected core (104
papers) under-counts true dependencies. Finally, this survey synthesises the field from
titles, abstracts, and citation structure rather than a close reading of every paper; the
per-paper claims, numbers, and attributions in
Sections~\ref{sec:genealogy}--\ref{sec:models} should therefore be read as a structural map
of the field rather than a verified account of each individual result.

\section{Conclusion}
Kazakh NLP is young, fast-growing, and structurally lopsided: rich in speech and
translation resources, thin where the language is morphologically hardest, and dependent
on a small set of datasets and global architectures. By mapping the field reproducibly and
ending with a concrete agenda, we hope to lower the barrier for the next wave of
contributors---especially those making a first contribution---and to make the morphology
and tokenization frontier, where the open questions are sharpest, easier to enter.

% ------------------------------------------------------------
% NOTE: author/year fields below need verification during the depth-read pass.
% Kazakh-resource authors (ISSAI et al.) are cited title-first to avoid misattribution.
% ------------------------------------------------------------
\begin{thebibliography}{99}

\bibitem{transformer} A.~Vaswani et al. Attention Is All You Need. \emph{NeurIPS}, 2017. arXiv:1706.03762.
\bibitem{bert} J.~Devlin et al. BERT: Pre-training of Deep Bidirectional Transformers. \emph{NAACL}, 2019. arXiv:1810.04805.
\bibitem{bpe} R.~Sennrich, B.~Haddow, A.~Birch. Neural Machine Translation of Rare Words with Subword Units. \emph{ACL}, 2016. arXiv:1508.07909.
\bibitem{xlmr} A.~Conneau et al. Unsupervised Cross-lingual Representation Learning at Scale (XLM-R). \emph{ACL}, 2020. arXiv:1911.02116.
\bibitem{bahdanau} D.~Bahdanau, K.~Cho, Y.~Bengio. Neural Machine Translation by Jointly Learning to Align and Translate. \emph{ICLR}, 2015. arXiv:1409.0473.
\bibitem{gpt3} T.~Brown et al. Language Models are Few-Shot Learners (GPT-3). \emph{NeurIPS}, 2020. arXiv:2005.14165.
\bibitem{rust} P.~Rust et al. How Good is Your Tokenizer? \emph{ACL}, 2021. arXiv:2012.15613.
\bibitem{arnett70} C.~Arnett et al. Evaluating Morphological Alignment of Tokenizers in 70 Languages. 2025. arXiv:2507.06378.
\bibitem{saleva} J.~S\"alev\"a, C.~Lignos. The Effectiveness of Morphology-aware Segmentation in Low-Resource NMT. 2021. arXiv:2103.11189.

\bibitem{ksc} Y.~Khassanov, S.~Mussakhojayeva, A.~Mirzakhmetov, A.~Adiyev, M.~Nurpeiissov, H.~A.~Varol. A Crowdsourced Open-Source Kazakh Speech Corpus and Initial Speech Recognition Baseline. 2020. arXiv:2009.10334.
\bibitem{kaznerd} KazNERD: Kazakh Named Entity Recognition Dataset. \emph{LREC}, 2022. arXiv:2111.13419.
\bibitem{kazakhtts} S.~Mussakhojayeva et al. KazakhTTS: An Open-Source Kazakh Text-to-Speech Synthesis Dataset. \emph{Interspeech}, 2021. arXiv:2104.08459.
\bibitem{kazakhtts2} KazakhTTS2: Extending the Open-Source Kazakh TTS Corpus. \emph{LREC}, 2022. arXiv:2201.05771.
\bibitem{ksc2} KSC2: An Industrial-Scale Open-Source Kazakh Speech Corpus. \emph{Interspeech}, 2022.
\bibitem{kazqad} KazQAD: Kazakh Open-Domain Question Answering Dataset. 2024. arXiv:2404.04487.
\bibitem{kazsandra} KazSAnDRA: Kazakh Sentiment Analysis Dataset of Reviews and Attitudes. 2024. arXiv:2403.19335.
\bibitem{dollms} A.~Maxutov, A.~Myrzakhmet, P.~Braslavski. Do LLMs Speak Kazakh? A Pilot Evaluation of Seven Models. \emph{SIGTURK}, 2024. ACL Anthology 2024.sigturk-1.8.
\bibitem{kazmmlu} KazMMLU: Evaluating Language Models on Kazakh, Russian, and Regional Knowledge. 2025. arXiv:2502.12829.
\bibitem{tumlu} TUMLU: A Unified Native Language Understanding Benchmark for Turkic Languages. 2025. arXiv:2502.11020.
\bibitem{qorgau} Qorgau: Evaluating LLM Safety in Kazakh-Russian Bilingual Contexts. 2025. arXiv:2502.13640.

\bibitem{sherkala} Sherkala-Chat: Building a State-of-the-Art LLM for Kazakh in a Moderately Resourced Setting. 2025. arXiv:2503.01493.
\bibitem{sozkz} SozKZ: Training Efficient Small Language Models for Kazakh from Scratch. 2026. arXiv:2603.20854.
\bibitem{kazbyte} KazByte: Adapting Qwen Models to Kazakh via a Byte-level Adapter. 2026. arXiv:2603.27859.

\bibitem{loreslm} Y.~Veitsman, M.~Hartmann. Recent Advancements and Challenges of Turkic Central Asian Language Processing. \emph{Proc. First Workshop on Language Models for Low-Resource Languages (LoResLM)}, 2025. arXiv:2407.05006.
\bibitem{turkicreview} P.~Muhetaer. Natural Language Processing and Speech Technologies for Central Asian Turkic Languages: A Review of Current Methods, Resources, and Challenges. \emph{Actual Problems of the Present}, 2025.

\end{thebibliography}

\end{document}
